\section*{Capítulo \ref{ch:g11:prob} --- Probabilidad y variables aleatorias}

\begin{solution}{exo:g11:prob:1}
modelo uniforme en $52$ tarjetas. Copas: $\frac{13}{52} = \frac14$. Reyes:
$\frac{4}{52} = \frac{1}{13}$. Corazón o rey:
$\frac{13 + 4 - 1}{52} = \frac{16}{52} = \frac{4}{13}$ (el rey de
corazones se cuenta una vez). Ni: $1 - \frac{4}{13} = \frac{9}{13}$.
\end{solution}

\begin{solution}{exo:g11:prob:2}
$ S = 6$: las parejas $(1,5), (2,4), (3,3), (4,2), (5,1)$, entonces
$\P(S = 6) = \frac{5}{36}$.

$ S \geq 10$: tres pares para $10$, dos para $11$, uno para $12$:
$\P(S \geq 10) = \frac{6}{36} = \frac16$.

$ S $ incluso: sumas $2, 4, 6, 8, 10, 12$ tener $1 + 3 + 5 + 5 + 3 + 1 = 18$
pares: \omterm{def:g10:proba:distribution}{probabilidad} $\frac{18}{36} = \frac12$.
\end{solution}

\begin{solution}{exo:g11:prob:3}
Las \omterm{def:g10:proba:distribution}{probabilidades} suman $1$: $ p = 1 - 0.3 - 0.2 - 0.4 = 0.1$. Entonces
\[
\E(X) = 0.3 \times (-1) + 0.2 \times 0 + 0.4 \times 2 + 0.1 \times 5
= -0.3 + 0.8 + 0.5 = 1 .
\]
\end{solution}

\begin{solution}{exo:g11:prob:4}
$\E(Y) = 2 \times 3 - 5 = 1$; $\V(Y) = 2^2 \times 4 = 16$;
$\sigma(Y) = 4$.
\end{solution}

\begin{solution}{exo:g11:prob:5}
Las ganancias son $0$, $2$, $5$ euros con \omterm{def:g10:proba:distribution}{probabilidades} $0.5$, $0.3$,
$0.2$; restando el $2$-apuesta en euros, $ G $ toma los valores $-2, 0, 3$:
\begin{center}
\begin{tabular}{c|ccc}
$ g $ & $-2$ & $0$ & $3$\\
\hline
$\P(G=g)$ & $0.5$ & $0.3$ & $0.2$
\end{tabular}
\end{center}
$\E(G) = -1 + 0 + 0.6 = -0.4$: el jugador pierde $40$ centavos por juego en
promedio --- desfavorable.
\end{solution}

\begin{solution}{exo:g11:prob:6}
Dibujo en orden: $\P(X = 2) = \frac35 \times \frac24 = \frac{3}{10}$
(rojo luego rojo); $\P(X = 0) = \frac25 \times \frac14 = \frac{1}{10}$
(azul luego azul); por eso
$\P(X = 1) = 1 - \frac{3}{10} - \frac{1}{10} = \frac{6}{10}$. Entonces
\[
\E(X) = 0 \times \frac1{10} + 1 \times \frac6{10} + 2 \times \frac3{10}
= \frac{12}{10} = 1.2 ,
\]
consistente con la intuición: $\frac35$ de los dos sorteos son rojos en promedio.
\end{solution}

\begin{solution}{exo:g11:prob:7}
$\E(G^2) = 0.5 \times 4 + 0.3 \times 0 + 0.2 \times 9 = 3.8$, entonces por el
fórmula de atajo
$\V(G) = 3.8 - (-0.4)^2 = 3.8 - 0.16 = 3.64$ y
$\sigma(G) \approx 1.9$ euros.
\end{solution}

\begin{solution}{exo:g11:prob:8}
El beneficio es $80$ menos el reclamo:
\begin{center}
\begin{tabular}{c|ccc}
$ b $ & $80$ & $-920$ & $-4920$\\
\hline
$\P(B=b)$ & $0.94$ & $0.05$ & $0.01$
\end{tabular}
\end{center}
$\E(B) = 75.2 - 46 - 49.2 = -20$: a este precio la empresa \emph{pierde}
$20$ euros por póliza de \omterm{def:g11:stat:mean}{media}. La prima es demasiado baja para los riesgos.
cubierto (tendría que exceder $100$ euros por una \omterm{def:g11:prob:expectation}{esperanza} positiva
beneficio).
\end{solution}

\begin{solution}{exo:g11:prob:9}
La puntuación aleatoria es $+1$ con \omterm{def:g10:proba:distribution}{probabilidad} $\frac14$ y $ x $ con
\omterm{def:g10:proba:distribution}{probabilidad} $\frac34$:
\[
\E = \frac14 + \frac{3x}{4} = 0 \iff x = -\frac13 .
\]
con la pena $-\frac13$, las conjeturas aleatorias no aportan nada en promedio.
\end{solution}

\begin{solution}{exo:g11:prob:10}
Dos monedas: coinciden (HH o TT) con \omterm{def:g10:proba:distribution}{probabilidad} $\frac24 \times 2 =
\frac12$. La ganancia del jugador es $+1$ o $-1$ con iguales \omterm{def:g10:proba:distribution}{probabilidades}:
\omterm{def:g11:prob:expectation}{esperanza} $0$ para ambos jugadores --- el juego es justo.

Tres monedas: todas coinciden (HHH o TTT) con \omterm{def:g10:proba:distribution}{probabilidad} $\frac{2}{8} =
\frac14$. El jugador recibe $2$ con \omterm{def:g10:proba:distribution}{probabilidad} $\frac14$ y paga $1$
con \omterm{def:g10:proba:distribution}{probabilidad} $\frac34$:
$\E = \frac24 - \frac34 = -\frac14$. El juego favorece a Ann por $25$ centavos
por ronda; sería justo si ella pagara $3$ euros por un triple partido.
\end{solution}

\begin{solution}{exo:g11:prob:11}
\emph{1.} El organizador recibe $1000 \times 5 = 5000$ euros y paga
fuera $2000 + 5 \times 200 + 20 \times 25 = 3500$ euros en total: el beneficio
es el \emph{constante} $1500$ (sin aleatoriedad una vez vendidos todos los boletos).
La ganancia de un solo jugador $ G $ (premio menos el $5$-billete en euros):
\begin{center}
\begin{tabular}{c|cccc}
$ g $ & $1995$ & $195$ & $20$ & $-5$\\
\hline\rule{0pt}{11pt}%
$\P(G=g)$ & $\dfrac{1}{1000}$ & $\dfrac{5}{1000}$ & $\dfrac{20}{1000}$
& $\dfrac{974}{1000}$
\end{tabular}
\end{center}

\emph{2.}
\[
\E(G) = \frac{1995 + 5 \times 195 + 20 \times 20 - 974 \times 5}{1000}
= \frac{1995 + 975 + 400 - 4870}{1000} = -1.5 .
\]
Cada jugador espera perder. $1.50$ euros; encima $1000$ jugadores que es
$1500$ euros --- exactamente el beneficio del organizador. El saldo de los libros:
lo que los jugadores pierden en promedio es precisamente lo que el organizador
recoge.
\end{solution}

\begin{solution}{pb:g11:prob:1}
\textbf{1.} $ G $ acepta $-2, -1, 0, 1, 2, 3$, cada uno con
\omterm{def:g10:proba:distribution}{probabilidad} $\frac16$: $\E(G) = \frac{-2 - 1 + 0 + 1 + 2 +
3}{6} = \frac12$ euro.

\textbf{2.} $\E(G^2) = \frac{4 + 1 + 0 + 1 + 4 + 9}{6} =
\frac{19}{6}$; $ V(G) = \frac{19}{6} - \frac14 = \frac{35}{12}
\approx 2.92$; $\sigma(G) \approx 1.71$.

\textbf{3.} $\E(2G - 1) = 2 \times \frac12 - 1 = 0$;
$ V(2G - 1) = 4\,V(G) = \frac{35}{3}$.

\textbf{4.} $\E(\text{sum}) = 7$: precio justo $7$ euros. En $8$,
la pérdida esperada es $1$ euros por partido.

\textbf{5.} Prima justa: $0.05 \times 800 = 40$ euros.
Cargando $60$: beneficio esperado $20$ euros por póliza.

\textbf{6.} Terrible $5:3$ recompensa el \emph{pasado} --- pero el
El bote es de quien gane. \emph{futuro}, y B podría
todavía gana. Darle todo a A considera una victoria probable como segura ---
La pequeña oportunidad de B vale algo.

\textbf{7.} Rondas a jugar: $3$; resultados (ganancia A/ganancia B por
ronda): AAA, AAB, ABA, ABB, BAA, BAB, BBA, BBB. El jugador B toma
la serie sólo ganando los tres: BBB solo. A gana el
serie en $7$ de $8$ casos: división justa
$64 \times \frac78 = 56$ pistolas para A, $8$ para b.

\textbf{8.} Jugar las rondas extra sin sentido cambia
victoria de nadie: una vez que A logra su sexta victoria, las rondas posteriores son
ceremonia. Pero fijando el horizonte exactamente $3$ rondas hace
todos $8$ resultados igualmente probables --- necesidades de conteo honesto
historias de la misma extensión y completarlas no cuesta nada.

\textbf{9.} En $5$--$5$: una ronda decisiva, $32$ cada. En
$5$--$4$: A obtiene $\frac{64 + 32}{2} = 48$. En $5$--$3$: el
la siguiente ronda conduce a $64$ (A lo gana) o al $5$--$4$
valor de la posición $48$: participación de A $\frac{64 + 48}{2} = 56$ ---
Fermat's $56 : 8$, re-derivado caminando hacia atrás.

\textbf{10.} una necesidad $2$, B necesita $3$: jugar $4$ imaginario
rondas ($16$ resultados). B toma la serie con $3$ o $4$
B-gana: $4 + 1 = 5$ resultados. Dividir:
$ B $: $64 \times \frac{5}{16} = 20$; $ A $: $44$.

\textbf{11.} El valor presente razonable de un futuro incierto
la recompensa es su \omterm{def:g11:prob:expectation}{esperanza} --- promedio ponderado por \omterm{def:g10:proba:distribution}{probabilidad} de
lo que pueda venir. Seguros (pregunta 5) y finanzas (precios
recompensas inciertas) ¿esta frase se convierte en industrias;
también lo son los casinos, desde el otro lado de la mesa.

\textbf{12.} El bote siempre termina en alguna parte: las dos acciones
son \omterm{def:g11:prob:rv}{variables aleatorias} sumando a la constante $64$, entonces por
linealidad $\E_A + \E_B = 64$ --- las acciones justas agotan el bote,
en cada puntuación.

\textbf{13.} Linealidad: $\E(X + Y) = 3.5 + 3.5 = 7$, sin mesa
--- cierto incluso para dados pegados. varianza de un dado:
$\E(X^2) - 3.5^2 = \frac{91}{6} - \frac{49}{4} = \frac{35}{12}$;
para dados independientes
$ V(X + Y) = \frac{35}{12} + \frac{35}{12} = \frac{35}{6}$.

\textbf{14.} Cada invitado fijo recupera su propio sombrero con
\omterm{def:g10:proba:distribution}{probabilidad} $\frac1n $ (su sombrero es uno entre $ n $). sumando
sobre el $ n $ huéspedes:
$\E(X) = n \times \frac1n = 1$ --- un invitado afortunado en promedio,
si el partido tiene dos sombreros o dos mil.

\textbf{15.} $ n = 2$: asignaciones: ambas correctas ($ X = 2$) o
intercambiado ($ X = 0$): $\E = 1$. $ n = 3$: las seis asignaciones dan
$ X = 3, 1, 1, 1, 0, 0$: $\E = \frac{6}{6} = 1$.

\textbf{16.} El cálculo añadió el \omterm{def:g11:prob:expectation}{esperanzas} de heredar del $ n $
indicadores individuales (``invitado $ i $ consiguió su sombrero"), y
la linealidad agrega \omterm{def:g11:prob:expectation}{esperanzas} de heredar \emph{incondicionalmente} ---
la dependencia arruinaría un producto o una varianza, pero nunca un
suma de \omterm{def:g11:prob:expectation}{esperanzas} de heredar. Ese es el superpoder de la linealidad: no
preguntas de independencia formuladas.

\textbf{17.} Ambas ganancias: $\E = 0.5 \times 10 - 5 = 0$ y
$0.005 \times 1000 - 5 = 0$. variaciones (de la ganancia):
R: $\E(G^2) = 0.5 \times 25 + 0.5 \times 25 = 25$;
B: $0.005 \times 995^2 + 0.995 \times 25 \approx 4\,975$:
$\sigma_A = 5$, $\sigma_B \approx 70.5$. La gente compra B: ¿qué?
ellos compran no es \omterm{def:g11:prob:expectation}{esperanza} (cero en ambos) pero el
\emph{dispersión} --- una pequeña posibilidad de obtener un derecho que cambie la vida
cola.

\textbf{18.} Primera cruz en el flip $ n $ tiene \omterm{def:g10:proba:distribution}{probabilidad}
$2^{-n}$ y paga $2^n $: cada $ n $ contribuye
$2^{-n} \times 2^n = 1$ hacia \omterm{def:g11:prob:expectation}{esperanza}, que por lo tanto es
$1 + 1 + 1 + \dots $: infinito. Sin embargo, ninguna persona en su sano juicio paga
$1\,000$ jugar una vez: con \omterm{def:g10:proba:distribution}{probabilidad} $\frac{15}{16}$ el
el pago es como \omterm{def:g10:functions:extrema}{máximo} $16$. \omterm{def:g11:prob:expectation}{esperanza} resume el largo plazo de
\emph{repetible} apuestas; un juego único en la vida con un comodín
la cola no tiene precio por su \omterm{def:g11:stat:mean}{media}.

\textbf{19.} \omterm{def:g11:prob:expectation}{Esperanzas} de heredar: asegurar $-20$; ir desnudo
$-0.01 \times 1000 = -10$. La estrategia desnuda gana en promedio
--- para alguien que pueda absorber la pérdida y repetir la apuesta
muchas veces. Una familia que no puede sobrevivir $-1\,000$ no es
en el negocio del promedio: vale la pena $10$ extra para eliminar el
varianza. Promedio de aseguradoras; los hogares compran certeza.

\textbf{20.} Nací dividiendo una olla entre dos impacientes.
jugadores; potenciado por la linealidad, que suma \omterm{def:g11:prob:expectation}{esperanzas} sin
pedir permiso a la dependencia; ciego al riesgo, que varianza
--- las loterías, el juego de San Petersburgo, la política familiar
--- debe medir en su lugar; y coronado el año próximo por la ley de
números grandes, lo que explica por qué los promedios siempre,
eventualmente, recoger.
\end{solution}
